% unit: noether.p07.sec.003
\subsection*{3.}
% unit: noether.p07.par.006
由刚刚证明的结果可以推出关于有理表示的结论。每一个有理绝对不变量都可以表示为两个整有理绝对不变量的商，这两个不变量不必互素；这是通过把分子和分母乘以分母关于该群的共轭即可直接看出的。因此，每一个有理绝对不变量都可以由 Galois 预解式 $\Phi(z,u)$ 的系数 $G_{\alpha\alpha_1\ldots\alpha_n}(x)$ 有理地表示。这个定理也可以不经由有限性定理而以多种方式容易地证明；它在 Weber II, \S\ 58 中已经表述。\footnote{然而，那里给出的证明并不可靠。它只表明公式 (7) 中的函数 $\Psi(t)$ 以不变量为系数，却没有表明这些不变量能由 $\Phi(t)$ 的系数有理表示。若在用 $\xi_k$ 表示 $x_i^{(k)}$ 时，不使用 Lagrange 插值公式，而使用一个已知的微分方法，则可避免这个缺口。这就是对恒等式 $\Phi(-\xi_k,u)=0$ 关于所有 $u$ 微分得到的关系：
\[
 \left[\frac{\partial\Phi}{\partial u_i}
      -x_i^{(k)}\frac{\partial\Phi}{\partial z}\right]_{z=-\xi_k}=0.
\]
相应地，对于每一个有理函数 $\omega(x)$，Weber 中的公式 (7) 和 (8) 应由下式代替：
\[
 \omega(x_1^{(k)},\ldots,x_n^{(k)})=
 \left.
 \omega\!\left(
   \frac{\partial\Phi/\partial u_1}{\partial\Phi/\partial z},\ldots,
   \frac{\partial\Phi/\partial u_n}{\partial\Phi/\partial z}
 \right)\right|_{z=-\xi_k},
\]
它只含有 $\Phi(z,u)$ 的系数；当 $\omega(x)$ 为不变量时，对所有 $k$ 求和即给出所要求的表示。}

% unit: noether.p07.sec.004
\subsection*{4.}
% unit: noether.p07.par.007
最后要指出，这里推出的结果已隐含于我的论文“有理函数的域与系统”之中；\footnote{Math. Ann. 76, p. 161 (1915).}更准确地说，1. 中给出的完整系统包含在定理 VI 和 VII 中，而独立于这个有限性定理的有理表示证明包含在定理 III 中。\footnote{对于相对不变量，定理 VIII 和 IX 包含了一个有限性证明，它同样建立在不同于通常证明的基础上，但也只是存在性证明；原因在于相对不变量并不构成一个域。}不过，在目前这个特殊情形中，证明过程和结果相对于一般理论都大为简化。我之所以想到把一般研究应用于有限群不变量，是由于与 E. Fischer 的谈话；2. 中给出的证明在内容上也来自他，3. 的注释亦然；在一般情形中，那里出现的完整系统并不能有理定义。

% unit: noether.p07.close.001
Erlangen, 1915 年 5 月。

\clearpage
% unit: noether.p08.title.001
